\section{Metodología}
\subsection{Participaciones y geometría composicional}
Sea $x_i=(x_{i1},\ldots,x_{iD})$ la composición media del departamento $i$, con $D=\Sectors$, $x_{ij}\geq0$ y $\sum_jx_{ij}=100$. El cierre implica dependencia entre las partes: aumentar una participación requiere que otra disminuya. Trabajamos con logaritmos de cocientes, siguiendo el enfoque composicional de Aitchison \cite{aitchison}.

El algoritmo admite reemplazo multiplicativo de ceros \cite{martin}. Si $z_i$ es el número de partes nulas y $\delta=\DeltaZero$ puntos de participación, la transformación es
\begin{equation}
x^*_{ij}=\begin{cases}\delta,&x_{ij}=0,\\x_{ij}(100-z_i\delta)/100,&x_{ij}>0.\end{cases}
\end{equation}
Se aplica después de normalizar la composición a 100. En los promedios efectivamente utilizados se observan \textbf{\Zeros\ ceros}; por tanto, el reemplazo no modifica los datos. Valores pequeños de extracción en Lambayeque y San Martín no son ceros y no se sustituyen. La sensibilidad a $\delta$ es aquí una comprobación sin perturbación efectiva, no evidencia de robustez ante ceros estructurales.

Las coordenadas CLR son
\begin{equation}
\operatorname{clr}(x^*_{ij})=\log x^*_{ij}-\frac{1}{D}\sum_{h=1}^{D}\log x^*_{ih}.
\end{equation}
Sus filas suman cero y viven en un subespacio de dimensión $D-1$. Calculamos también coordenadas ILR con una base ortonormal de Helmert \cite{egozcue}. La distancia euclidiana entre vectores ILR coincide con la distancia de Aitchison entre composiciones.

Aplicamos PCA a CLR centrado, sin estandarizar columnas. Los primeros dos componentes explican \PcaCumulative\% de la varianza (\PcaOne\% y \PcaTwo\%); tres conservan \PcaThree\%. El primer eje se relaciona con extracción y el segundo contrasta manufactura con agropecuario, según sus cargas. El signo de cada componente es convencional. El mapa no superpone flechas de cargas en una escala arbitraria: las cargas se reportan por separado en el apéndice.

\subsection{Agrupamiento y validación}
Comparamos Ward, K-Means, PAM y mezclas gaussianas para K=2,3,4,5 sobre CLR+PCA2, CLR+PCA3 e ILR completo: \SpecCount\ configuraciones. K-Means utiliza 20 inicializaciones; PAM busca intercambios desde medoides iniciales; GMM utiliza cinco inicializaciones y se evalúa además mediante BIC. PAM es una búsqueda local, sin garantía de óptimo global. La semilla fija facilita reproducir los ejercicios.

Evaluamos silueta, Calinski--Harabasz y Davies--Bouldin. La silueta depende de la representación: un valor mayor tras reducir dimensión no demuestra que los componentes descartados sean ruido económico irrelevante. La partición principal de cuatro perfiles se utiliza por interpretabilidad y estabilidad relativa, manteniendo alternativas como sensibilidad.

Calculamos Gap con \GapReplicas\ referencias uniformes en la caja de las coordenadas PCA observadas \cite{gap}. Aplicamos la regla del primer K que cumple $\operatorname{Gap}(K)\geq\operatorname{Gap}(K+1)-s_{K+1}$. La estabilidad se aproxima con \StabilityReplicas\ submuestras del \SubsamplePct\% \textbf{sin reemplazo}: cada grupo original, restringido a las unidades seleccionadas, se compara con el grupo remuestreado de mayor Jaccard. Promediamos por grupo y luego por réplica. Este procedimiento es submuestreo, una alternativa discutida por Hennig \cite{hennig}; no es bootstrap ni un intervalo de confianza.

Los nombres K=4 se asignan por participación extractiva media ordenada. Son etiquetas descriptivas. En particular, \textit{diversificado/agro} incluye Lima, Lambayeque y San Martín; comparten baja extracción, sin que los tres sean economías predominantemente agrarias. Tampoco se deduce diversificación de una sola proporción minera; el HHI sectorial se exporta como descriptor complementario.

Distinguimos el corte Ward K=2 de una agregación económica de K=4. El primero separa \WardSmallCount\ y \WardLargeCount\ departamentos. La segunda agrupa los dos perfiles de mayor extracción en un bloque minero de \MacroMiningCount\ unidades y reúne los demás en otro de \MacroOtherCount. Se utiliza en M2 y no equivale al corte Ward de mayor silueta.

\subsection{Volatilidad, crecimiento e índice de nivel}
El crecimiento anual y el índice respecto a 2019 son
\begin{equation}
g_{it}=100\left(\frac{VAB_{it}}{VAB_{i,t-1}}-1\right),\qquad I_{it}=100\frac{VAB_{it}}{VAB_{i,2019}}.
\end{equation}
Una caída de 12\% seguida de un rebote de 9\% produce $100(0.88)(1.09)=95.92$. Es un ejemplo aritmético, no una trayectoria estimada. Las métricas informan dimensiones distintas.

La desviación estándar por perfil utiliza las observaciones departamento-año de $g_{it}$. Incluye heterogeneidad entre unidades, variación temporal, pandemia y rebote; no equivale a volatilidad idiosincrática tras controlar choques comunes. Reportamos Levene centrado en la mediana y Fligner--Killeen como diagnósticos descriptivos. Sus valores $p$ convencionales no ajustan la dependencia serial o transversal del panel.

\subsection{Modelos de efectos fijos}
Para los modelos categóricos estimamos
\begin{equation}\begin{split}
y_{it}={}&\alpha_i+\beta_C C_t+\beta_P P_t\\
&+\sum_{k\ne r}\gamma_k C_tG_{ik}+\sum_{k\ne r}\theta_k P_tG_{ik}+u_{it},
\end{split}\end{equation}
donde $C_t=1[t=2020]$, $P_t=1[2021\leq t\leq2023]$ y $r$ es diversificado/agro. La dependiente es $g_{it}$ en M1 y M4 e $I_{it}$ en M5. Los efectos fijos absorben los niveles constantes departamentales y los indicadores de grupo, que no varían en el tiempo.

M2 sustituye los cuatro perfiles por la agregación económica de dos grupos. M3 y M6 sustituyen las categorías por $z_i$, participación extractiva previa estandarizada con la media y desviación estándar muestral de los 24 departamentos. M3 usa crecimiento y M6 índice. M4 replica M1 excluyendo Apurímac 2016. La tabla \ref{tab:modelos} resume las especificaciones.

\begin{table}[htbp]\centering\small\begin{tabular}{llrrr}
\toprule
Modelo & Especificación & Periodo & Obs. & $R^2_w$ \\
\midrule
M1 & Tasas, cuatro perfiles & 2008--2023 & 384 & 0.145 \\
M2 & Tasas, dos macrogrupos & 2008--2023 & 384 & 0.141 \\
M3 & Tasas, minería continua & 2008--2023 & 384 & 0.142 \\
M4 & Tasas, sin Apurímac 2016 & 2008--2023 & 383 & 0.277 \\
M5 & Índice, cuatro perfiles & 2017--2023 & 168 & 0.355 \\
M6 & Índice, minería continua & 2017--2023 & 168 & 0.345 \\
\bottomrule
\end{tabular}

\caption{Especificaciones y muestras}\label{tab:modelos}
\nota{Todas incluyen efectos fijos departamentales. M4 retira una observación. El $R^2$ within no permite ordenar modelos con diferente dependiente y muestra.}\end{table}

Los coeficientes $\theta_k$ comparan el cambio entre \PostStart--\PostEnd\ y el periodo omitido con el cambio del grupo de referencia. En M5, el periodo omitido es \LevelStart--\PreEnd; en M1 comprende los años previos a 2020 de su ventana más larga. No es una brecha de nivel de un solo año ni la distancia frente a una trayectoria contrafactual sin pandemia. La base 2019 iguala el índice de todos los departamentos en ese año, pero no iguala sus trayectorias previas.

Los errores estándar se agrupan por departamento mediante \texttt{PanelOLS}; los intervalos se obtienen con \texttt{conf\_int} y la distribución t de la versión registrada del estimador \cite{linearmodels}. Se conserva el ajuste por regresores (\texttt{debiased=True}) y la configuración predeterminada sin ajuste adicional por número de grupos. Hay \Departments\ grupos de inferencia y solo \EnclaveCount\ departamentos enclave. No se incluyen efectos fijos anuales completos: los indicadores comunes recogen solo los regímenes definidos. El diseño no garantiza tendencias paralelas, ausencia de confusión temporal o asignación exógena a una estructura productiva.
\FloatBarrier
